\documentclass[12pt]{article}
\usepackage[utf8]{inputenc}
\usepackage[T1]{fontenc}
\usepackage{hyperref}
\title{Upper Baram Kenyah Language Resources for K-12 Curriculum Development}
\author{Compiled by Claude Sonnet 5 (claude-sonnet-5) for Andr\'as Kornai}
\date{}
\begin{document}
\maketitle

\section{Introduction}
Upper Baram Kenyah (ISO 639-3 ubm) belongs to the Kenyah branch of the North Sarawakan languages, which are Austronesian. It is spoken in the upper Baram area of northeastern Sarawak, Malaysia. The identification line in the task gives the family as ``Bookkeeping''. This is evidently a data error, and the Austronesian classification comes from the search results. Smith (2017) places the Lepo' varieties (Lepo' Gah, Lepo' Tau, Lepo' Sawa and others) within a Highland Kenyah subgroup. The searches did not settle how those varieties relate to Upper Baram Kenyah as a separate language. Joshua Project, citing Ethnologue 2016, lists about 900 speakers for a Kenyah group, but its Indonesia-framed page may not apply to Sarawak. The searches found no school materials.

\section{Primary Educational Resources}
No primers, textbooks or school materials were found.

\section{Linguistic and Reference Materials}
\begin{itemize}
\item Wikipedia, North Sarawakan languages (classification): \url{https://en.wikipedia.org/wiki/North_Sarawakan_languages}
\item Library of Congress subject authority, Upper Baram Kenyah language (alternate names): \url{https://id.loc.gov/authorities/subjects/sh2007007017.html}
\item Borneo Dictionary, Kenyah Lebo' entries (Lepo' Gah speech from Long San; example page): \url{https://borneodictionary.com/kenyah-lebo-vo/long-selatong}
\end{itemize}

\section{Religious and Community Resources}
\begin{itemize}
\item Joshua Project, Upper Baram Kenyah in Malaysia: \url{https://joshuaproject.net/index.php/people_groups/12625/MY}
\item Borneo Post article on Kenyah communities of the upper Baram: \url{https://www.theborneopost.com/?p=1317770}
\end{itemize}

\section{Digital and Audio Resources}
No audio, video or corpus resources were found.

\section{Curriculum Development Context}
The searches found little. A curriculum would need new documentation work, and the Lepo' varieties may differ enough to need a choice of standard. Malaysian schools use Malay as the medium of instruction. The searches did not show whether any Kenyah mother-tongue programme exists.
\end{document}